\subsection{Angles and the Calculation of Their Measure}

Right angles have already been used as Euclidean constructions: they are what
make the axes perpendicular and what make the triangle in the distance formula a
right triangle. We now make the general language of angles explicit and define
how to assign numerical measures to angles.

\begin{definitionbox}
Let $r_1$ and $r_2$ be two rays with the same initial point $O$. The rays,
together with a choice of one of the two regions between them, determine an
angle $\alpha$ with vertex $O$.
\end{definitionbox}

The angle is the chosen geometric region. It is neither the small arc drawn
near the vertex nor a number. The arc is only a graphical mark indicating
which region has been selected.

\begin{center}
\begin{tikzpicture}[scale=1.0]
    \coordinate (O) at (0,0);
    \coordinate (A) at (3.4,0);
    \coordinate (B) at (1.7,2.8);
    \fill[blue!8] (O) -- (A) arc[start angle=0,end angle=58.7,radius=3.4] -- cycle;
    \draw[axis] (O) -- (A) node[right] {$r_1$};
    \draw[axis] (O) -- (B) node[above] {$r_2$};
    \draw[subset] (0.9,0) arc[start angle=0,end angle=58.7,radius=0.9];
    \node[subset label] at (0.95,0.48) {$\alpha$};
    \fill (O) circle (2pt) node[below left] {$O$};
\end{tikzpicture}
\end{center}

\subsubsection*{Why circles give a consistent measure}

Draw a circle centred at $O$. The rays cut from the circle an arc corresponding
to the chosen angle.

\begin{theorembox}
For every Euclidean circle, the ratio of circumference $C$ to diameter $D$ is
the same constant:
\[
\frac{C}{D}=\pi.
\]
This constant is called $\pi$. Since $D=2r$, every circle of radius $r$ has
circumference
\[
\boxed{C=2\pi r}.
\]
\end{theorembox}

The constancy of $C/D$ follows from similarity: enlarging a circle by a factor
$\lambda$ multiplies every length, including its diameter, circumference, and
corresponding arc lengths, by $\lambda$.

If the same two rays meet two circles of radii $r_1$ and $r_2$ and cut arc
lengths $s_1$ and $s_2$, then
\[
\frac{s_1}{s_2}=\frac{r_1}{r_2},
\]
and therefore
\[
\frac{s_1}{r_1}=\frac{s_2}{r_2}.
\]
Thus the quotient $s/r$ depends only on the angle, not on the chosen circle.

\begin{definitionbox}
Let $s$ be the length of the selected arc on a circle of radius $r$ centred at
the vertex. The radian measure of $\alpha$ is
\[
\boxed{m_{\mathrm{rad}}(\alpha)=\frac{s}{r}}.
\]
\end{definitionbox}

For a full turn, $s=2\pi r$, so
\[
m_{\mathrm{rad}}(\text{full turn})=2\pi.
\]
A half-turn has radian measure $\pi$, and a quarter-turn has radian measure
$\pi/2$.

Degrees are a second numerical scale. A full turn is assigned $360^\circ$, so
\[
2\pi\text{ radians}=360^\circ
\]
and
\[
\boxed{
m_{\deg}(\alpha)
=
\frac{180^\circ}{\pi}\,m_{\mathrm{rad}}(\alpha)
}.
\]

\subsubsection*{The law of cosines is inherited from Euclidean geometry}

Directly measuring an arc is often difficult. We therefore use a theorem that
belongs to the Euclidean geometry with which we started.

Choose points
\[
A\in r_1\setminus\{O\},
\qquad
B\in r_2\setminus\{O\}.
\]
The points form the Euclidean triangle $OAB$. Before writing the formula, we
identify the angle and the three side lengths occurring in it.

\begin{center}
\begin{tikzpicture}[scale=0.92]
    \coordinate (O) at (0,0);
    \coordinate (A) at (4.4,0);
    \coordinate (B) at (1.55,3.0);

    \fill[blue!6] (O) -- (1.0,0)
        arc[start angle=0,end angle=62.7,radius=1.0] -- cycle;

    \draw[subset] (O) -- (A)
        node[midway,below=5pt] {$d_E(O,A)$};
    \draw[subset] (O) -- (B)
        node[midway,left=7pt] {$d_E(O,B)$};
    \draw[subset] (A) -- (B)
        node[midway,above right=3pt] {$d_E(A,B)$};

    \draw[subset] (0.78,0)
        arc[start angle=0,end angle=62.7,radius=0.78];
    \node[subset label] at (0.88,0.42) {$\alpha$};

    \fill (O) circle (2.1pt) node[below left] {$O$};
    \fill (A) circle (2.1pt) node[below right] {$A$};
    \fill (B) circle (2.1pt) node[above] {$B$};

    \node[align=left,anchor=west] at (5.0,2.35) {
        adjacent to $\alpha$:\\
        $d_E(O,A)$ and $d_E(O,B)$
    };
    \node[align=left,anchor=west] at (5.0,0.65) {
        opposite $\alpha$:\\
        $d_E(A,B)$
    };
\end{tikzpicture}
\end{center}

The side $AB$ is opposite the angle $\alpha$, while $OA$ and $OB$ are adjacent
to it. With these symbols fixed by the diagram, the classical law of cosines,
proved geometrically without Cartesian coordinates, states that
\[
d_E(A,B)^2
=
d_E(O,A)^2+d_E(O,B)^2
-2d_E(O,A)d_E(O,B)
\cos\bigl(m_{\mathrm{rad}}(\alpha)\bigr).
\]

\begin{remarkbox}
The law of cosines is not created or justified by the coordinate system in this
lecture. It is an earlier theorem of Euclidean geometry. Coordinates enter only
when we calculate the three side lengths
\[
d_E(O,A),\qquad d_E(O,B),\qquad d_E(A,B)
\]
from the coordinates of $O$, $A$, and $B$.
\end{remarkbox}

Solving the Euclidean law of cosines for the angle measure gives
\begin{theorembox}
\[
\boxed{
m_{\mathrm{rad}}(\alpha)
=
\arccos\left(
\frac{
d_E(O,A)^2+d_E(O,B)^2-d_E(A,B)^2
}{
2d_E(O,A)d_E(O,B)
}
\right)
}.
\]
The distance formula supplies the three numerical lengths needed in this
expression.
\end{theorembox}

The logical direction is therefore
\[
\begin{aligned}
&\text{Euclidean geometry and the law of cosines}\\
&\qquad\longrightarrow\text{Cartesian coordinates}\\
&\qquad\longrightarrow\text{numerical side lengths}\\
&\qquad\longrightarrow\text{calculation of the angle measure}.
\end{aligned}
\]

\begin{examplebox}
Let
\[
O=(0,0),\qquad A=(1,0),\qquad B=(1,1).
\]
The coordinate distance formula gives
\[
d_E(O,A)=1,
\qquad
d_E(O,B)=\sqrt2,
\qquad
d_E(A,B)=1.
\]
The Euclidean law of cosines then gives
\[
\cos\bigl(m_{\mathrm{rad}}(\alpha)\bigr)
=
\frac{1^2+(\sqrt2)^2-1^2}{2\cdot1\cdot\sqrt2}
=
\frac{1}{\sqrt2}.
\]
Hence
\[
m_{\mathrm{rad}}(\alpha)=\frac{\pi}{4},
\qquad
m_{\deg}(\alpha)=45^\circ.
\]
The angle is the geometric region determined by the rays $OA$ and $OB$; the
numbers $\pi/4$ and $45^\circ$ are two measures of that angle.
\end{examplebox}

\begin{definitionbox}
A \emph{right angle} is the constructible angle formed by two perpendicular
rays. Equivalently, two perpendicular lines meet so that the four angles around
their intersection are congruent; each of these angles is called a right angle.
\end{definitionbox}

This name was already used earlier in the construction of the axes and in the
Pythagorean triangle for the distance formula. The new information obtained now
is numerical: if $\alpha$ is a right angle, then
\[
m_{\mathrm{rad}}(\alpha)=\frac{\pi}{2},
\qquad
m_{\deg}(\alpha)=90^\circ.
\]
A right angle is not either of these numbers; it is a geometric angle having
these measures in the two measurement systems.

\begin{summarybox}
The construction proceeds in the following order:
\begin{enumerate}
    \item begin with a Euclidean plane containing no chosen coordinates;
    \item construct one straight line;
    \item choose a point $O$ on that line;
    \item construct through $O$ the perpendicular line, and only then name the
          two lines as coordinate axes and $O$ as their origin;
    \item choose positive directions on both axes;
    \item choose one common unit length and transfer it with a compass in both
          positive and negative directions to construct the integer marks;
    \item repeatedly bisect intervals to construct the dense family of dyadic
          marks;
    \item use Euclidean continuity to identify every axis position with a real
          number;
    \item assign to every point its unique pair of orthogonal projections;
    \item derive the coordinate distance formula from the resulting right
          triangle and the Pythagorean theorem;
    \item only after this construction introduce the numerical measurement of
          general angles.
\end{enumerate}
\end{summarybox}
